\documentclass[12pt,a4paper]{article}

% --- Idioma y codificación ---
\usepackage[utf8]{inputenc}
\usepackage[T1]{fontenc}
\usepackage[spanish,es-noshorthands]{babel}
\usepackage{lmodern}

% --- Matemáticas ---
\usepackage{amsmath,amssymb,amsfonts}

% --- Márgenes y espaciado ---
\usepackage[margin=2.5cm]{geometry}
\usepackage{enumitem}

\title{Banco de Ejercicios de \textit{FUNctions}:\\ Hoja de Práctica}
\author{Edwin Ernesto Rodríguez Pérez\\ Danny Mauricio Trujillo\\ María Erika González Ortez}
\date{}

\begin{document}

\maketitle

\noindent\textbf{Instrucciones:} Resuelve cada ejercicio a lápiz, mostrando el procedimiento completo y justificando cada respuesta.

\vspace{1em}

\begin{enumerate}[label=\textbf{\arabic*.}, itemsep=2.5em]

% 1
\item El panel primario de energía muestra el conjunto de entrada $A=\{1, 2, 3\}$ y el conjunto de salida $B=\{4, 5\}$. La relación actual de transmisión es $R = \{(1,4), (2,4), (3,5)\}$. Evalúa las conexiones y determina si la corriente fluirá de forma válida (¿Es una función formal?).

% 2
\item ¡Alerta de sobrecarga! El nodo $1 \in A$ está enviando señal a dos puertos distintos en $B$: $R = \{(1,4), (1,5), (2,4), (3,5)\}$. Desconecta el cable excedente para restablecer la unicidad y convertir la relación en una función.

% 3
\item El circuito está incompleto: $R = \{(1,4), (3,5)\}$ sobre $A=\{1, 2, 3\}$ y $B=\{4, 5\}$. El nodo $2 \in A$ no tiene señal. Conecta el nodo $2$ a cualquier puerto de $B$ para completar la condición de existencia.

% 4
\item El escáner analiza $f(x) = 2x - 1$, con codominio real y dominio real máximo. Determina su dominio y rango mediante las proyecciones sobre los ejes $X$ e $Y$.

% 5
\item El sensor sigue $f(x)=\sqrt{x-2}$, tomando la raíz principal y codominio real. Determina el dominio real máximo y el rango. Indica si sus extremos están incluidos.

% 6
\item La señal sigue $f(x)=5/(x-3)$, con codominio real. Determina su dominio real máximo e identifica el valor de $x$ excluido.

% 7
\item Sintoniza la antena alineando la trayectoria rectilínea $f(x) = mx + b$. Determina $m$ y $b$ para conectar los puntos de control $(0,4)$ y $(2,0)$. Responde ambos valores y escribe la ecuación resultante.

% 8
\item Construye la trayectoria $f(x)=(x-h)^2+k$, con dominio real, para que su vértice sea $(3,2)$. Ajusta $h$ y $k$; el coeficiente $a$ permanece fijo en $1$. Escribe la ecuación de la parábola que construiste.

% 9
\item El escáner refleja una intensidad de salida $y = 8$ sobre la función cuadrática $f(x) = x^2 - 1$. Encuentra TODAS las preimágenes $x$ que producen esa señal.

% 10
\item Analiza el generador espacial $f(x)=2^x$. Determina su dominio real máximo, su rango y la ecuación de su asíntota horizontal.

% 11
\item Refleja los puntos $(0,1)$ y $(1,2)$ de $f(x)=2^x$ sobre la recta $y=x$ para construir su inversa $g(x)=\log_2(x)$. Indica los dos puntos reflejados, el dominio, el rango y la asíntota vertical de $g$.

% 12
\item El núcleo emite un impulso de salida $y = 81$ bajo la regla $f(x) = 3^x$. Transforma la ecuación a su forma logarítmica equivalente para hallar el valor exacto del exponente $x$.

% 13
\item El filtro de datos principal opera con la función $f(x) = \log_5(x - 4)$. Resuelve la restricción analítica del argumento para fijar los límites del escáner y la asíntota de falla.

% 14
\item El sistema de diagnóstico analiza la llegada de una señal de nivel $y = -2$ en dos módulos distintos:
  \begin{enumerate}[label=\arabic*), leftmargin=*, itemsep=0.2em]
    \item \textit{Módulo Exponencial:} $f(x)=2^x$; resolver $f(x)=-2$ en los reales.
    \item \textit{Módulo Lineal:} $g(x)=3x+1$; resolver $g(x)=-2$ en los reales.
  \end{enumerate}
  Determina el conjunto de preimágenes $S_1$ y $S_2$ para cada módulo.

% 15
\item El núcleo principal del satélite funciona con el sistema por tramos:
  \[
    f(x) =
    \begin{cases}
      x + 2 & \text{si } x < 0 \\[4pt]
      (x - 3)^2 + 2 & \text{si } x \ge 0
    \end{cases}
  \]
  Evalúa la respuesta del sistema para $f(-2)$ y $f(0)$, y determina el dominio total de la función.

\end{enumerate}

\end{document}
